\documentclass[../main]{subfiles}
\begin{document}

\chapter*{Preface to New Typesetting}
\addcontentsline{toc}{chapter}{Preface to New Typesetting}

Typesetting credits go to: Aareyan Manzoor, Carl Sun, George Coote, RokettoJanpu, Prakhar Agarwal, Yohan Wittgenstein, Shuayb Mohammed, Arpit Mittal, Manan Jain. %add name here

Here is a link to a dyslexic friendly version: \url{https://aareyanmanzoor.github.io/assets/matsumura-CA-dyslexic.pdf}

We changed some notation. Maps between spec was changed from $^a\phi$ to $\SpecInduced{\phi}$. $I$-adic completion was changed from $A^\ast$ to $\completion{A}$.

Some comments by editors were made in footnotes.

We added a bibliography, the original book cited references in the text without having a dedicated bibliography.

We added citations and references with hyperlinks. References to e.g. theorems/paragraphs in the book are in blue, while citations to the bibliography is in red. The bibliography also has URLs now, for easy access. Some of the books in the bibliography  had newer editions, so we went with those.

This revision restores mathematical text, hypotheses, differentials, diagrams and references checked against the 1980 second edition. It also restores the separate sixteen-entry Index of Symbols, with page references to this typesetting, and corrects the EGA fascicles and Northcott bibliography record.

Corrections of errors already present in the original are distinguished from transcription restorations by explicit editorial notes. These include the completion chapter reference, the primary-decomposition inclusion, the zero-module depth convention, the assertion about isomorphisms of extensions, differential-module subscripts, the compactness argument, references in the excellence chapter, the $p$-degree and differential-basis statements, formal etaleness, the Kunz proof's multiplication and height formula, and Ratliff's notation. The generic-fibre argument in \S\ref{sec:42} makes the original's implicit localization explicit. In the Kunz conjecture the missing hypothesis is $R^p\subseteq S\subseteq R$, as stated by Kimura--Niitsuma; the previous editorial note had reversed it.

The paragraph markers in \S\ref{sec:20} and \S\ref{sec:27} and the symbolic references in \S\ref{sec:31} have been restored. The existing consecutive labels (14.I) and (14.J), corresponding to the original's (14.K) and (14.L), are retained; their internal references are consistent. The notation modernizations described above are unchanged.

\end{document}